\documentclass{article}
\usepackage[a4paper, total={6in, 8in}]{geometry}
\usepackage{verbatim}
\usepackage{hyperref}

\title{VibeVolts Documentation}
\author{Jules}
\date{\today}

\begin{document}

\maketitle

\tableofcontents

\section{constants.py}

This file defines global constants that are used throughout the VibeVolts project. These constants represent array indices and physical properties, which helps to improve code readability and maintainability by avoiding the use of hardcoded "magic numbers". The constants defined in this file include:

\begin{itemize}
    \item \textbf{Radii in Meters:} Defines the radii of celestial bodies like the Earth and Moon.
    \item \textbf{Detector Array Indices:} Specifies the column indices for various detector properties in a NumPy array, such as aperture size, pixel size, quantum efficiency, and exclusion angles.
    \item \textbf{Orbital Elements Array Indices:} Details the indices for orbital parameters like the semi-major axis, eccentricity, and inclination.
    \item \textbf{Pointing State Array Indices:} Outlines the indices for the pointing state, including the total number of points and the current position in the sequence.
\end{itemize}

\section{constellation.py}

This file is responsible for creating constellations of satellites. It contains functions to generate orbital elements for satellites in specific configurations, such as a geostationary constellation.

\subsection{geos(sim\_data, n, fov)}
\begin{itemize}
    \item \textbf{Purpose:} Creates `n` equally spaced satellites in a geostationary orbit (GEO) and adds them to the simulation data. It calculates the required pointing grid based on the satellite's field of view (FOV) to ensure full sky coverage. It also initializes the orbital elements and pointing states for each satellite.
    \item \textbf{Arguments:}
    \begin{itemize}
        \item `sim\_data`: The main simulation data dictionary.
        \item `n`: The number of satellites to create.
        \item `fov`: The diameter of the field of view of the satellite in radians.
    \end{itemize}
    \item \textbf{Returns:} `None`. It modifies `sim\_data` in place.
\end{itemize}

\section{generate\_log\_spherical\_points.py}

This file provides a function to generate 3D points that are logarithmically spaced radially and uniformly distributed angularly. This is useful for creating a point cloud where the density of points decreases with distance from the origin.

\subsection{generate\_log\_spherical\_points(num\_points, inner\_radius, outer\_radius, object\_size\_m=1.0, seed=None)}
\begin{itemize}
    \item \textbf{Purpose:} Generates a specified number of 3D points with logarithmic radial spacing and uniform angular distribution on spherical shells using the Fibonacci lattice method. Each point is also assigned a specific size.
    \item \textbf{Arguments:}
    \begin{itemize}
        \item `num\_points`: The total number of points to generate.
        \item `inner\_radius`: The minimum distance from the origin.
        \item `outer\_radius`: The maximum distance from the origin.
        \item `object\_size\_m`: The size in meters to be associated with each point.
        \item `seed`: An optional integer to seed the random number generator for reproducible results.
    \end{itemize}
    \item \textbf{Returns:} A tuple containing two NumPy arrays: one for the 3D Cartesian coordinates of the points, and one for the size of each object.
\end{itemize}

\section{lambertian.py}

This file contains a function to calculate the brightness of a Lambertian sphere, which is a sphere that reflects light diffusely.

\subsection{lambertiansphere(vec\_from\_sphere\_to\_light, vec\_from\_sphere\_to\_observer, albedo, radius)}
\begin{itemize}
    \item \textbf{Purpose:} Calculates the effective brightness of a Lambertian sphere. The function determines the apparent brightness based on the angle between the light source and the observer, the sphere's albedo (reflectivity), and its size.
    \item \textbf{Arguments:}
    \begin{itemize}
        \item `vec\_from\_sphere\_to\_light`: A 3-element NumPy array representing the direction vector from the sphere to the light source.
        \item `vec\_from\_sphere\_to\_observer`: A 3-element NumPy array representing the direction vector from the sphere to the observer.
        \item `albedo`: The fraction of incident light that is reflected (0.0 to 1.0).
        \item `radius`: The radius of the sphere in meters.
    \end{itemize}
    \item \textbf{Returns:} The effective brightness cross-section in square meters, which is proportional to the total light reflected towards the observer.
\end{itemize}

\section{observatories.py}

This file is responsible for adding observatory data structures to the main simulation data.

\subsection{add\_observatories(sim\_data, num\_observatories)}
\begin{itemize}
    \item \textbf{Purpose:} Initializes and adds data structures for a specified number of observatories to the simulation data dictionary. This function sets up zero-filled NumPy arrays for the observatories' positions, velocities, accelerations, pointing vectors, and detector properties.
    \item \textbf{Arguments:}
    \begin{itemize}
        \item `sim\_data`: The main simulation data dictionary.
        \item `num\_observatories`: The number of observatories to add.
    \end{itemize}
    \item \textbf{Returns:} `None`. It modifies `sim\_data` in place.
\end{itemize}

\section{plotting\_3d.py}

This file provides functionality for creating 3D scatter plots of object positions, which is essential for visualizing the spatial distribution of objects in the simulation.

\subsection{plot\_3d\_scatter(positions, title, plot\_time, labels=None, marker\_size=1, trace\_name='Points')}
\begin{itemize}
    \item \textbf{Purpose:} Creates a 3D scatter plot of object positions using Plotly. The plot includes a representation of the Earth, its equator, the North Pole, and a specific location (El Segundo). It is designed to visualize the positions of satellites or other objects in a 3D space relative to the Earth.
    \item \textbf{Arguments:}
    \begin{itemize}
        \item `positions`: A NumPy array of shape (n, 3) containing the Cartesian coordinates of the points to plot.
        \item `title`: The title of the plot.
        \item `plot\_time`: The `datetime` for which the plot is generated, used for coordinate transformations.
        \item `labels`: An optional list of strings to label each point.
        \item `marker\_size`: The size of the markers for the points.
        \item `trace\_name`: The name of the trace for the points in the legend.
    \end{itemize}
    \item \textbf{Returns:} A Plotly `go.Figure` object representing the 3D plot.
\end{itemize}

\section{plotting\_vectors.py}

This file is dedicated to creating 3D visualizations of satellites along with their pointing vectors, which is crucial for understanding the orientation and direction of satellites in the simulation.

\subsection{plot\_pointing\_vectors(data\_struct, title, plot\_time)}
\begin{itemize}
    \item \textbf{Purpose:} Generates a 3D plot that displays satellites as markers and their corresponding pointing vectors as lines. The plot includes a 3D model of the Earth to provide a frame of reference. This visualization is useful for verifying that satellites are pointing in the expected directions.
    \item \textbf{Arguments:}
    \begin{itemize}
        \item `data\_struct`: The main simulation data dictionary, which contains the positions and pointing vectors of the satellites.
        \item `title`: The title of the plot.
        \item `plot\_time`: The `datetime` for which the plot is generated.
    \end{itemize}
    \item \textbf{Returns:} A Plotly `go.Figure` object representing the 3D plot.
\end{itemize}

\section{pointing.py}

This file manages satellite pointing behavior, including updating pointing directions, generating pointing grids, and handling cases where satellites have no visibility.

\subsection{pointing\_place\_update(data\_struct)}
\begin{itemize}
    \item \textbf{Purpose:} Increments the pointing index for all satellites, allowing them to cycle through a predefined pointing grid. It wraps the index back to 0 when it reaches the end of the sequence.
    \item \textbf{Arguments:}
    \begin{itemize}
        \item `data\_struct`: The main simulation data dictionary.
    \end{itemize}
    \item \textbf{Returns:} `None`. It modifies `data\_struct` in place.
\end{itemize}

\subsection{jerk(data\_struct, satellite\_number)}
\begin{itemize}
    \item \textbf{Purpose:} Randomly adjusts the pointing vector of a specific satellite by a fixed angle (0.3 radians). This is used to reorient satellites that are "blind" (i.e., cannot see any target points).
    \item \textbf{Arguments:}
    \begin{itemize}
        \item `data\_struct`: The main simulation data dictionary.
        \item `satellite\_number`: The index of the satellite to modify.
    \end{itemize}
    \item \textbf{Returns:} The modified `data\_struct`.
\end{itemize}

\subsection{generate\_pointing\_sphere(data\_struct, n\_points)}
\begin{itemize}
    \item \textbf{Purpose:} Creates a uniform grid of pointing vectors on a sphere (a "pointing sphere") using the `pointing_vectors` function and stores it in the `data_struct` if it doesn't already exist for the given number of points.
    \item \textbf{Arguments:}
    \begin{itemize}
        \item `data\_struct`: The main simulation data dictionary.
        \item `n\_points`: The number of points to generate for the pointing sphere.
    \end{itemize}
    \item \textbf{Returns:} `None`. It modifies `data\_struct` in place.
\end{itemize}

\subsection{update\_satellite\_pointing(data\_struct)}
\begin{itemize}
    \item \textbf{Purpose:} Updates each satellite's pointing vector to the next position in its assigned pointing grid based on its current pointing state.
    \item \textbf{Arguments:}
    \begin{itemize}
        \item `data\_struct`: The main simulation data dictionary.
    \end{itemize}
    \item \textbf{Returns:} `None`. It modifies `data\_struct` in place.
\end{itemize}

\subsection{find\_and\_jerk\_blind\_satellites(data\_struct)}
\begin{itemize}
    \item \textbf{Purpose:} Identifies satellites that cannot see any fixed points (based on the visibility table) and applies the `jerk` function to them to change their orientation.
    \item \textbf{Arguments:}
    \begin{itemize}
        \item `data\_struct`: The main simulation data dictionary.
    \end{itemize}
    \item \textbf{Returns:} The modified `data\_struct`.
\end{itemize}

\section{pointing\_vectors.py}

This file provides functions for generating and visualizing equally spaced points on a unit sphere, which are used as pointing vectors for satellites.

\subsection{pointing\_vectors(n)}
\begin{itemize}
    \item \textbf{Purpose:} Generates a specified number of equally spaced points on a unit sphere using the Fibonacci lattice algorithm. This creates a uniform distribution of vectors that can be used for systematic sky scanning.
    \item \textbf{Arguments:}
    \begin{itemize}
        \item `n`: The number of points to generate.
    \end{itemize}
    \item \textbf{Returns:} A NumPy array of shape (n, 3) containing the Cartesian coordinates of the points on the unit sphere.
\end{itemize}

\subsection{plot\_vectors\_on\_sphere(vectors, title)}
\begin{itemize}
    \item \textbf{Purpose:} Creates a 3D plot to visualize a set of vectors on a unit sphere. This is useful for verifying that the pointing vectors generated by `pointing_vectors` are distributed as expected.
    \item \textbf{Arguments:}
    \begin{itemize}
        \item `vectors`: A NumPy array of shape (n, 3) representing the vectors to plot.
        \item `title`: The title of the plot.
    \end{itemize}
    \item \textbf{Returns:} A Plotly `go.Figure` object representing the 3D plot.
\end{itemize}

\section{propagation.py}

This file handles the orbital propagation of satellites and celestial bodies. It includes functions to read satellite data from Two-Line Element (TLE) files, update the positions of the Sun and Moon, and propagate satellite orbits forward in time.

\subsection{add\_satellites\_from\_tle(sim\_data, tle\_file\_path, sat\_category)}
\begin{itemize}
    \item \textbf{Purpose:} Initializes a category of satellites in the simulation data by reading their orbital elements from a TLE file. It sets up the necessary data structures for position, velocity, and pointing, but does not calculate their initial positions.
    \item \textbf{Arguments:}
    \begin{itemize}
        \item `sim\_data`: The main simulation data dictionary.
        \item `tle\_file\_path`: The path to the TLE file.
        \item `sat\_category`: The key for this satellite category (e.g., 'satellites').
    \end{itemize}
    \item \textbf{Returns:} `None`. It modifies `sim\_data` in place.
\end{itemize}

\subsection{celestial\_update(data\_struct, time\_date)}
\begin{itemize}
    \item \textbf{Purpose:} Calculates and updates the positions of the Sun and Moon in the Geocentric Celestial Reference System (GCRS) for a given time.
    \item \textbf{Arguments:}
    \begin{itemize}
        \item `data\_struct`: The main simulation data dictionary.
        \item `time\_date`: The timezone-aware `datetime` for which to calculate the positions.
    \end{itemize}
    \item \textbf{Returns:} The modified `data\_struct`.
\end{itemize}

\subsection{readtle(tle\_file\_path)}
\begin{itemize}
    \item \textbf{Purpose:} Reads a TLE file and extracts the orbital elements and epoch for each satellite. It returns the elements in a canonical order within a NumPy array.
    \item \textbf{Arguments:}
    \begin{itemize}
        \item `tle\_file\_path`: The path to the TLE file.
    \end{itemize}
    \item \textbf{Returns:} A tuple containing a NumPy array of orbital elements and a list of `datetime` epochs.
\end{itemize}

\subsection{propagate\_satellites\_new(data\_struct, time\_date, sat\_category=None)}
\begin{itemize}
    \item \textbf{Purpose:} Updates the positions of satellites based on their orbital elements by propagating them from their epoch to a new time. It also sets the default pointing vector for each satellite to be radially outward from the Earth's center.
    \item \textbf{Arguments:}
    \begin{itemize}
        \item `data\_struct`: The main simulation data dictionary.
        \item `time\_date`: The `datetime` to which the satellites should be propagated.
        \item `sat\_category`: An optional string to specify a single satellite category to propagate.
    \end{itemize}
    \item \textbf{Returns:} The modified `data\_struct`.
\end{itemize}

\section{radiometry\_calcs.py}

This file provides a suite of functions for performing radiometric calculations. These include converting between linear flux ratios and astronomical magnitudes, calculating blackbody radiation properties using Planck's law and the Stefan-Boltzmann law, and generating plots of blackbody spectra.

\subsection{mag(x)}
\begin{itemize}
    \item \textbf{Purpose:} Calculates an astronomical magnitude from a linear flux ratio.
    \item \textbf{Arguments:}
    \begin{itemize}
        \item `x`: The positive linear flux ratio.
    \end{itemize}
    \item \textbf{Returns:} The corresponding magnitude value.
\end{itemize}

\subsection{amag(x)}
\begin{itemize}
    \item \textbf{Purpose:} Calculates the linear flux ratio from an astronomical magnitude. It is the inverse of the `mag` function.
    \item \textbf{Arguments:}
    \begin{itemize}
        \item `x`: The magnitude value.
    \end{itemize}
    \item \textbf{Returns:} The corresponding linear flux ratio.
\end{itemize}

\subsection{\_planck\_law(wav\_m, temp\_k)}
\begin{itemize}
    \item \textbf{Purpose:} A helper function that implements Planck's law to calculate the spectral radiance of a blackbody at a specific wavelength and temperature.
    \item \textbf{Arguments:}
    \begin{itemize}
        \item `wav\_m`: The wavelength in meters.
        \item `temp\_k`: The temperature in Kelvin.
    \end{itemize}
    \item \textbf{Returns:} The spectral radiance in W / (m\textasciicircum 2 * sr * m).
\end{itemize}

\subsection{blackbody\_flux(temperature, lambda\_short, lambda\_long)}
\begin{itemize}
    \item \textbf{Purpose:} Numerically integrates the spectral radiance of a blackbody over a given wavelength band to compute the total flux within that band.
    \item \textbf{Arguments:}
    \begin{itemize}
        \item `temperature`: The temperature of the blackbody in Kelvin.
        \item `lambda\_short`: The short wavelength of the band in microns.
        \item `lambda\_long`: The long wavelength of the band in microns.
    \end{itemize}
    \item \textbf{Returns:} The integrated spectral radiance in Watts / (m\textasciicircum 2 * steradian).
\end{itemize}

\subsection{stefan\_boltzmann\_law(temperature)}
\begin{itemize}
    \item \textbf{Purpose:} Calculates the total power radiated per unit area by a blackbody across all wavelengths, according to the Stefan-Boltzmann law.
    \item \textbf{Arguments:}
    \begin{itemize}
        \item `temperature`: The temperature of the blackbody in Kelvin.
    \end{itemize}
    \item \textbf{Returns:} The total radiated power per unit area in W / m\textasciicircum 2.
\end{itemize}

\subsection{plot\_blackbody\_spectrum(temperature)}
\begin{itemize}
    \item \textbf{Purpose:} Generates a log-log plot of a blackbody's spectral radiance over the mid-infrared range (0.5 to 30 microns).
    \item \textbf{Arguments:}
    \begin{itemize}
        \item `temperature`: The temperature of the blackbody in Kelvin.
    \end{itemize}
    \item \textbf{Returns:} `None`. Displays a Plotly figure.
\end{itemize}

\subsection{plot\_blackbody\_spectrum\_visible\_nir(temperature)}
\begin{itemize}
    \item \textbf{Purpose:} Generates a plot of a blackbody's spectral radiance over the visible and near-infrared range (0.1 to 1 micron).
    \item \textbf{Arguments:}
    \begin{itemize}
        \item `temperature`: The temperature of the blackbody in Kelvin.
    \end{itemize}
    \item \textbf{Returns:} `None`. Displays a Plotly figure.
\end{itemize}

\section{radiometry\_data.py}

This file serves as a data repository for radiometric information, primarily containing physical constants and a comprehensive dictionary of standard astronomical filter data. The data is essential for calculating the brightness of celestial objects as observed through different filters.

\begin{itemize}
    \item \textbf{Physical Constants:} Defines constants such as the Astronomical Unit (`AU_M`) and the Sun's radius (`RSUN_M`) in meters.
    \item \textbf{FILTER\_DATA Dictionary:} A dictionary where each key is a string representing a standard astronomical filter (e.g., 'U', 'B', 'V', 'J', 'H', 'K', 'g', 'r', 'i', 'z', 'L', 'M', 'N', and several JWST filters). The value for each key is another dictionary containing the following data for that filter:
    \begin{itemize}
        \item `sun`: The apparent magnitude of the Sun in that filter.
        \item `sky`: The sky brightness in magnitudes per square arcsecond for a dark ground-based site.
        \item `space`: The background brightness for a space-based observatory.
        \item `central_wavelength`: The central wavelength of the filter in nanometers.
        \item `bandwidth`: The bandwidth of the filter in nanometers.
        \item `zero_point`: The photon flux for a 0-magnitude object in photons/sec/m\textasciicircum 2.
    \end{itemize}
\end{itemize}

\section{simulation.py}

This file provides functions to initialize and set up the core components of the space simulation environment. It includes functions for creating the main simulation data structure, adding celestial bodies, and generating fixed reference points.

\subsection{create\_empty\_simulation(start\_time, delta\_time=60.0)}
\begin{itemize}
    \item \textbf{Purpose:} Initializes a minimal, empty data structure for a space simulation. It sets up the basic dictionary that will hold all simulation data, including the start time and time step.
    \item \textbf{Arguments:}
    \begin{itemize}
        \item `start\_time`: The starting time and date of the simulation (must be a timezone-aware `datetime` object).
        \item `delta\_time`: The time step for the simulation in seconds.
    \end{itemize}
    \item \textbf{Returns:} A dictionary representing the basic simulation state.
\end{itemize}

\subsection{add\_celestial\_bodies(sim\_data)}
\begin{itemize}
    \item \textbf{Purpose:} Adds the data structures for celestial bodies (specifically the Sun and Moon) to the simulation data dictionary.
    \item \textbf{Arguments:}
    \begin{itemize}
        \item `sim\_data`: The main simulation data dictionary.
    \end{itemize}
    \item \textbf{Returns:} `None`. It modifies `sim\_data` in place.
\end{itemize}

\subsection{add\_fixed\_points(sim\_data, num\_points=100)}
\begin{itemize}
    \item \textbf{Purpose:} Adds a structure for fixed reference points in the Geocentric Celestial Reference System (GCRS) frame. The points are generated with a logarithmic spherical distribution.
    \item \textbf{Arguments:}
    \begin{itemize}
        \item `sim\_data`: The main simulation data dictionary.
        \item `num\_points`: The number of fixed points to generate.
    \end{itemize}
    \item \textbf{Returns:} `None`. It modifies `sim\_data` in place.
\end{itemize}

\section{visibility.py}

This file contains functions to determine whether a satellite's line of sight is obstructed by celestial bodies like the Sun, Moon, or Earth. It is crucial for simulating realistic observational constraints.

\subsection{solarexclusion(data\_struct)}
\begin{itemize}
    \item \textbf{Purpose:} Performs a vectorized calculation to determine which satellites are affected by solar exclusion. It computes the angle between each satellite's pointing vector and the vector to the Sun, flagging those that fall within the specified exclusion angle.
    \item \textbf{Arguments:}
    \begin{itemize}
        \item `data\_struct`: The main simulation data dictionary.
    \end{itemize}
    \item \textbf{Returns:} A tuple containing two NumPy arrays: an `exclusion_vector` (1 if excluded, 0 otherwise) and an `angle_vector` (the calculated angle in radians for each satellite).
\end{itemize}

\subsection{exclusion(data\_struct, satellite\_index, print\_debug=False)}
\begin{itemize}
    \item \textbf{Purpose:} Checks if a single satellite's pointing vector is excluded by the Sun, Moon, or Earth. It considers the apparent size of the Moon and Earth when checking for exclusion.
    \item \textbf{Arguments:}
    \begin{itemize}
        \item `data\_struct`: The main simulation data dictionary.
        \item `satellite\_index`: The index of the satellite to check.
        \item `print\_debug`: If `True`, it prints detailed debug information.
    \end{itemize}
    \item \textbf{Returns:} `0` if the satellite's view is excluded, `1` otherwise.
\end{itemize}

\subsection{update\_visibility\_table(data\_struct, print\_debug\_for\_sat=None)}
\begin{itemize}
    \item \textbf{Purpose:} Updates the visibility table, which is a matrix indicating whether each satellite can see each fixed point. For each satellite-point pair, it sets the satellite's pointing vector towards the point and then calls the `exclusion` function to check for obstructions.
    \item \textbf{Arguments:}
    \begin{itemize}
        \item `data\_struct`: The main simulation data dictionary.
        \item `print\_debug\_for\_sat`: An optional satellite index for which to enable detailed debug printing in the `exclusion` function.
    \end{itemize}
    \item \textbf{Returns:} `None`. It modifies `data\_struct` in place.
\end{itemize}

\section{Demo Summaries}

This section provides a summary of the various demo files included in the VibeVolts project. Each demo is designed to illustrate a specific feature or capability of the simulation framework.

\begin{itemize}
    \item \textbf{all\_demos.py:} This script serves as a master runner for all other demo files. It imports the individual demo functions and executes them sequentially. It can either display the generated plots inline or save them to a single HTML file named \texttt{all\_demo\_plots.html}.
    \item \textbf{demo\_common.py:} This is not a standalone demo but a utility file used by other demos. It provides a function, `initialize_standard_simulation`, which sets up a standard simulation environment. This includes creating the main data structure, loading a predefined set of satellites from TLEs (Low Earth Orbit, Geostationary, and Highly Elliptical Orbit), adding celestial bodies (Sun and Moon), and generating fixed reference points. This ensures that demos start with a consistent and fully populated simulation state.
    \item \textbf{demo1.py:} This demo initializes a standard simulation and generates a 3D plot showing the initial positions of all satellites relative to the Earth at the simulation start time.
    \item \textbf{demo2.py:} This demo propagates the standard satellite constellation forward in time by one hour and then generates a 3D plot of their new positions. This demonstrates the orbital propagation functionality.
    \item \textbf{demo3.py:} This demo propagates the satellite constellation for 24 hours, one hour at a time, and generates a 3D plot showing the orbital tracks of all satellites over the full 24-hour period.
    \item \textbf{demo4.py:} This demo focuses on a single satellite (ISS) and plots its ground track on a 2D world map over a 24-hour period.
    \item \textbf{demo\_constellation.py:} This demo showcases the creation of a custom satellite constellation. It initializes an empty simulation, adds 10 geostationary satellites, and then generates a 3D plot to visualize their positions.
    \item \textbf{demo\_exclusion\_debug\_print.py:} This demo is designed to test and verify the line-of-sight exclusion logic. It initializes a standard simulation, updates the visibility table, and then runs the exclusion check for a specific satellite with debug printing enabled to show the detailed calculations for Sun, Moon, and Earth exclusion.
    \item \textbf{demo\_exclusion\_table.py:} This demo visualizes the satellite visibility for a set of fixed points. It initializes a simulation, updates the visibility table, and then creates a heatmap plot of the table, showing which satellites can see which fixed points.
    \item \textbf{demo\_fixedpoints.py:} This demo visualizes the distribution of the fixed reference points used in the simulation. It generates and plots these points in 3D space to confirm their logarithmic spherical distribution.
    \item \textbf{demo\_lambertian.py:} This demo illustrates the `lambertiansphere` function. It calculates and prints the effective brightness of a sphere (representing the Moon) as seen from another position (representing a satellite) under illumination from the Sun.
    \item \textbf{demo\_pointing\_plot.py:} This demo visualizes the pointing vectors of the satellites in the standard simulation. It generates a 3D plot showing each satellite and a vector indicating its current pointing direction.
    \item \textbf{demo\_pointing\_sequence.py:} This demo demonstrates how a satellite's pointing can be systematically updated. It shows a single satellite cycling through a predefined sequence of 25 pointing directions and plots the final pointing vector.
    \item \textbf{demo\_pointing\_vectors.py:} This demo visualizes a uniformly distributed set of pointing vectors on a sphere. It generates 1000 points using the Fibonacci lattice method and plots them on a unit sphere.
    \item \textbf{demo\_sky\_scan.py:} This demo simulates a sky survey. It initializes a constellation of 10 GEO satellites, each with a predefined pointing grid, and simulates them scanning the sky. It then generates a plot showing the number of visible fixed points over time.
    \item \textbf{show\_geo\_search.py:} This demo visualizes the process of a geostationary satellite searching for objects. It simulates a single GEO satellite scanning the sky and plots the number of new objects it discovers over time, demonstrating how the "jerk" function helps it find new targets when its view is empty.
\end{itemize}

\section{Code Listing: constants.py}
\begin{verbatim}
# --- Global Constants for Array Indices ---
# These constants define column indices for numpy arrays, making the
# code more readable and preventing errors from using "magic numbers".

# -- Radii in Meters --
EARTH_RADIUS = 6378137.0
MOON_RADIUS = 1737400.0

# -- Detector Array Indices --
DETECTOR_APERTURE_IDX = 0      # Aperture size in meters
DETECTOR_PIXEL_SIZE_IDX = 1    # Pixel size in radians
DETECTOR_QE_IDX = 2            # Quantum efficiency as a fraction (0.0 to 1.0)
DETECTOR_PIXELS_IDX = 3        # Total number of pixels in the detector (count)
DETECTOR_SOLAR_EXCL_IDX = 4    # Solar exclusion angle in radians
DETECTOR_LUNAR_EXCL_IDX = 5    # Lunar exclusion angle in radians
DETECTOR_EARTH_EXCL_IDX = 6    # Earth exclusion angle (above the limb) in radians

# -- Orbital Elements Array Indices --
ORBITAL_A_IDX = 0              # Semi-major axis in meters
ORBITAL_E_IDX = 1              # Eccentricity (dimensionless)
ORBITAL_I_IDX = 2              # Inclination in radians
ORBITAL_RAAN_IDX = 3           # Right Ascension of the Ascending Node in radians
ORBITAL_ARGP_IDX = 4           # Argument of Perigee in radians
ORBITAL_M_IDX = 5              # Mean Anomaly in radians

# -- Pointing State Array Indices --
POINTING_COUNT_IDX = 0         # Number of points in the pointing grid
POINTING_PLACE_IDX = 1         # Current index in the pointing sequence
\end{verbatim}

\section{Code Listing: constellation.py}
\begin{verbatim}
import numpy as np
from datetime import datetime
from typing import Dict, Any
import math
import random

from constants import (
    ORBITAL_A_IDX, ORBITAL_E_IDX, ORBITAL_I_IDX,
    ORBITAL_RAAN_IDX, ORBITAL_ARGP_IDX, ORBITAL_M_IDX,
    POINTING_COUNT_IDX, POINTING_PLACE_IDX
)
from pointing import generate_pointing_sphere

def geos(sim_data, n,  fov) -> None:
    """
    Creates n equally spaced satellites in GEO and adds them to the 'satellites' group in the
    simulation.

    Args:
        sim_data: The main simulation data dictionary.
        n: The number of satellites to create.
        fov: The diameter of the field of view of the satellite in radians.
    """
    # Calculate solid angle
    theta = fov / 2
    solid_angle = 2 * np.pi * (1 - np.cos(theta))

    # Calculate grid_points - blow things up by 0.25 for overlap
    grid_points = int(4 * np.pi / solid_angle * 1.25)

    # Generate and store the pointing sphere and place in ['pointing_sphers'][n]
    generate_pointing_sphere(sim_data, grid_points)

    orbital_elements_list = []
    epochs_list = []
    pointing_state_list = []

    # Geostationary orbit semi-major axis in meters
    a = 42164000.0

    # Create a set of elements evenly spaced around the equator in
    # orbital_elemens_list
    for i in range(n):
        elements = np.zeros(6)
        elements[ORBITAL_A_IDX] = a
        elements[ORBITAL_E_IDX] = 0.0
        elements[ORBITAL_I_IDX] = 0.0
        elements[ORBITAL_RAAN_IDX] = i * 2 * np.pi / n
        elements[ORBITAL_ARGP_IDX] = 0.0
        elements[ORBITAL_M_IDX] = 0.0
        orbital_elements_list.append(elements)
        epochs_list.append(sim_data['start_time'])

   # And these are the pointing state of the satellite. Make the position random.

        pointing_state = np.zeros(2, dtype=int)
        pointing_state[POINTING_COUNT_IDX] = grid_points
        pointing_state[POINTING_PLACE_IDX] = random.randint(0,grid_points-1)

        pointing_state_list.append(pointing_state)

    orbital_elements = np.array(orbital_elements_list, dtype=float)
    pointing_state_array = np.array(pointing_state_list, dtype=int)

    if 'satellites' not in sim_data:
        sim_data['counts']['satellites'] = n
        sim_data['satellites'] = {
            'position': np.zeros((n, 3), dtype=float),
            'velocity': np.zeros((n, 3), dtype=float),
            'acceleration': np.zeros((n, 3), dtype=float),
            'orbital_elements': orbital_elements,
            'epochs': epochs_list,
            'pointing': np.zeros((n, 3), dtype=float),
            'pointing_state': pointing_state_array,
            'detector': np.zeros((n, 7), dtype=float),
        }
    else:
        # Append to existing satellites
        sim_data['counts']['satellites'] += n
        sim_data['satellites']['position'] = np.vstack([sim_data['satellites']['position'],
            np.zeros((n, 3), dtype=float)])
        sim_data['satellites']['velocity'] = np.vstack([sim_data['satellites']['velocity'],
            np.zeros((n, 3), dtype=float)])
        sim_data['satellites']['acceleration'] = np.vstack([sim_data['satellites']['acceleration'],
            np.zeros((n, 3), dtype=float)])
        sim_data['satellites']['orbital_elements'] = np.vstack([
            sim_data['satellites']['orbital_elements'], orbital_elements])
        sim_data['satellites']['epochs'].extend(epochs_list)
        sim_data['satellites']['pointing'] = np.vstack([sim_data['satellites']['pointing'],
            np.zeros((n, 3), dtype=float)])
        sim_data['satellites']['pointing_state'] = np.vstack([
            sim_data['satellites']['pointing_state'], pointing_state_array])
        sim_data['satellites']['detector'] = np.vstack([sim_data['satellites']['detector'],
            np.zeros((n, 7), dtype=float)])
\end{verbatim}

\section{Code Listing: generate\_log\_spherical\_points.py}
\begin{verbatim}
import numpy as np
from datetime import datetime, timezone
from plotting_3d import plot_3d_scatter


def generate_log_spherical_points(
    num_points: int,
    inner_radius: float,
    outer_radius: float,
    object_size_m: float = 1.0,
    seed: int = None
) -> tuple[np.ndarray, np.ndarray]:
    """
    Generates 3D points with logarithmic radial and uniform angular distribution.

    This function creates a point cloud where point distances from the origin are
    logarithmically spaced. On any given spherical shell, points are distributed
    uniformly using the Fibonacci lattice method. Each point is associated with a
    specified object size.

    Args:
        num_points: The total number of points to generate.
        inner_radius: The minimum distance from the origin (must be positive).
        outer_radius: The maximum distance from the origin (must be >= inner_radius).
        object_size_m: The size in meters to be associated with each point.
                       Defaults to 1.0.
        seed: An optional integer to seed the random number generator for
              reproducible shuffling.

    Returns:
        A tuple containing:
        - A NumPy array of shape (num_points, 3) for the Cartesian coordinates.
        - A NumPy array of shape (num_points,) for the object size in meters.
    """
    # Input validation
    if not isinstance(num_points, int) or num_points <= 0:
        raise ValueError("num_points must be a positive integer.")
    if not (isinstance(inner_radius, (int, float)) and inner_radius > 0):
        raise ValueError("inner_radius must be a positive number.")
    if not (isinstance(outer_radius, (int, float)) and outer_radius >= inner_radius):
        raise ValueError("outer_radius must be a positive number and >= inner_radius.")

    # --- 1. Generate uniformly distributed unit vectors (Fibonacci Lattice) ---
    # Create an array of indices for each point
    indices = np.arange(0, num_points, dtype=float) + 0.5

    # Uniformly distribute points along the z-axis (cos(phi))
    z = 1 - 2 * indices / num_points

    # Calculate the radius in the xy-plane for each point
    radius_xy = np.sqrt(1 - z**2)

    # Calculate the azimuthal angle using the golden angle
    golden_angle = np.pi * (3. - np.sqrt(5.))
    theta = golden_angle * indices

    # Convert to Cartesian coordinates for the unit sphere
    x = radius_xy * np.cos(theta)
    y = radius_xy * np.sin(theta)

    # Stack into an (N, 3) array of unit vectors
    unit_vectors = np.stack([x, y, z], axis=1)

    # --- 2. Generate logarithmically spaced radii ---
    # The start and stop parameters for logspace are exponents
    start_exp = np.log10(inner_radius)
    stop_exp = np.log10(outer_radius)
    radii = np.logspace(start_exp, stop_exp, num_points)

    # --- 3. Shuffle radii to decouple radius from spherical position ---
    # This prevents the smallest radii from always mapping to points near the
    # north pole, creating a more uniform visual distribution of densities.
    rng = np.random.default_rng(seed)
    rng.shuffle(radii)

    # --- 4. Scale unit vectors by radii using broadcasting ---
    # Reshape radii to (N, 1) to broadcast with (N, 3) unit_vectors
    points = unit_vectors * radii[:, np.newaxis]

    # --- 5. Create an array for the object sizes ---
    sizes = np.full(num_points, object_size_m, dtype=float)

    return points, sizes

if __name__ == '__main__':
    # --- Demo of Point Generation and Visualization ---
    # This provides a consistent demonstration entry point, matching vibevolts_demo.py
    NUM_POINTS = 100
    # Using the same radii as the main simulation for consistency
    INNER_RADIUS = 2000000
    OUTER_RADIUS = 84328000
    OBJECT_SIZE = 1.0

    print(f"Generating {NUM_POINTS} points from radius {INNER_RADIUS} to {OUTER_RADIUS}...")
    generated_points, _ = generate_log_spherical_points(
        num_points=NUM_POINTS,
        inner_radius=INNER_RADIUS,
        outer_radius=OUTER_RADIUS,
        object_size_m=OBJECT_SIZE,
        seed=42
    )
    print("Point generation complete.\n")

    # Use the unified visualization function
    plot_3d_scatter(
        positions=generated_points,
        title="Fixed Points Distribution (from generate_log_spherical_points.py)",
        plot_time=datetime.now(timezone.utc),
        labels=[f"Point {i}" for i in range(len(generated_points))],
        marker_size=1,
        trace_name='Fixed Points'
    )
\end{verbatim}

\section{Code Listing: lambertian.py}
\begin{verbatim}
import numpy as np

def lambertiansphere(
    vec_from_sphere_to_light: np.ndarray,
    vec_from_sphere_to_observer: np.ndarray,
    albedo: float,
    radius: float
) -> float:
    """
    Calculates the effective brightness of a
    Lambertian sphere.

    This function determines the apparent brightness of
    a diffusely reflecting sphere based on the angle
    between the light source and the observer, the
    sphere's albedo (reflectivity), and its size.

    Args:
        vec_from_sphere_to_light: A 3-element NumPy
            array representing the direction vector from
            the sphere to the light source.
        vec_from_sphere_to_observer: A 3-element NumPy
            array representing the direction vector from
            the sphere to the observer.
        albedo: The fraction of incident light that is
            reflected (0.0 to 1.0).
        radius: The radius of the sphere in meters.

    Returns:
        The effective brightness cross-section in
        square meters. This value is proportional to
        the total light reflected towards the observer.
    """
    if not 0.0 <= albedo <= 1.0:
        raise ValueError("Albedo must be between 0.0 and 1.0.")
    if radius < 0:
        raise ValueError("Radius cannot be negative.")

    norm_light = np.linalg.norm(vec_from_sphere_to_light)
    norm_observer = np.linalg.norm(vec_from_sphere_to_observer)

    if norm_light == 0 or norm_observer == 0:
        raise ValueError("Input vectors cannot have zero length.")

    unit_vec_light = vec_from_sphere_to_light / norm_light
    unit_vec_observer = vec_from_sphere_to_observer / norm_observer

    cos_alpha = np.dot(unit_vec_light, unit_vec_observer)
    cos_alpha = np.clip(cos_alpha, -1.0, 1.0)
    alpha = np.arccos(cos_alpha)

    term1 = np.sin(alpha)
    term2 = (np.pi - alpha) * np.cos(alpha)
    phase_function_value = (2 / (3 * np.pi)) * (term1 + term2)

    cross_sectional_area = np.pi * (radius ** 2)

    effective_brightness = (
        albedo *
        cross_sectional_area *
        phase_function_value
    )

    return effective_brightness
\end{verbatim}

\section{Code Listing: observatories.py}
\begin{verbatim}
import numpy as np
from typing import Dict, Any

def add_observatories(sim_data: Dict[str, Any], num_observatories: int) -> None:
    """
    Adds observatory data structures to the simulation data.

    Args:
        sim_data: The main simulation data dictionary.
        num_observatories: The number of observatories to add.
    """
    if not isinstance(num_observatories, int) or num_observatories < 0:
        raise ValueError("num_observatories must be a non-negative integer.")

    sim_data['counts']['observatories'] = num_observatories
    sim_data['observatories'] = {
        'position': np.zeros((num_observatories, 3), dtype=float),
        'velocity': np.zeros((num_observatories, 3), dtype=float),
        'acceleration': np.zeros((num_observatories, 3), dtype=float),
        'pointing': np.zeros((num_observatories, 3), dtype=float),
        'detector': np.zeros((num_observatories, 7), dtype=float),
    }
\end{verbatim}

\section{Code Listing: plotting\_3d.py}
\begin{verbatim}
import numpy as np
from datetime import datetime
from typing import List, Optional

import plotly.graph_objects as go
from astropy.time import Time
from astropy.coordinates import GCRS, EarthLocation
import astropy.units as u

EARTH_RADIUS = 6378137.0

def plot_3d_scatter(
    positions: np.ndarray,
    title: str,
    plot_time: datetime,
    labels: Optional[List[str]] = None,
    marker_size: int = 1,
    trace_name: str = 'Points'
) -> go.Figure:
    """
    Creates a 3D plot of object positions.
    """
    if positions.ndim != 2 or positions.shape[1] != 3:
        raise ValueError("positions array must have shape (n, 3).")

    fig = go.Figure()

    fig.add_trace(go.Scatter3d(
        x=positions[:, 0], y=positions[:, 1], z=positions[:, 2],
        mode='markers',
        marker=dict(
            size=marker_size,
            color=np.arange(len(positions)),
            colorscale='Viridis',
            opacity=0.8
        ),
        text=labels,
        hoverinfo='text' if labels else 'none',
        name=trace_name
    ))

    u_sphere = np.linspace(0, 2 * np.pi, 100)
    v_sphere = np.linspace(0, np.pi, 100)
    x_earth = EARTH_RADIUS * np.outer(np.cos(u_sphere), np.sin(v_sphere))
    y_earth = EARTH_RADIUS * np.outer(np.sin(u_sphere), np.sin(v_sphere))
    z_earth = EARTH_RADIUS * np.outer(np.ones(np.size(u_sphere)), np.cos(v_sphere))
    fig.add_trace(go.Surface(
        x=x_earth, y=y_earth, z=z_earth,
        colorscale='Blues', showscale=False, opacity=0.5, name='Earth'
    ))

    theta = np.linspace(0, 2 * np.pi, 100)
    x_eq = EARTH_RADIUS * np.cos(theta)
    y_eq = EARTH_RADIUS * np.sin(theta)
    z_eq = np.zeros_like(theta)
    fig.add_trace(go.Scatter3d(
        x=x_eq, y=y_eq, z=z_eq, mode='lines',
        line=dict(color='green', width=3), name='Equator'
    ))

    fig.add_trace(go.Scatter3d(
        x=[0], y=[0], z=[EARTH_RADIUS * 1.1], mode='text',
        text=['N'], textfont=dict(size=15, color='red'), name='North Pole'
    ))

    lat_es, lon_es = 33.92 * u.deg, -118.42 * u.deg
    el_segundo_loc = EarthLocation.from_geodetic(lon=lon_es, lat=lat_es)
    itrs_coords = el_segundo_loc.get_itrs(obstime=Time(plot_time))
    gcrs_coords = itrs_coords.transform_to(GCRS(obstime=Time(plot_time)))
    es_pos = gcrs_coords.cartesian.xyz.to(u.m).value * 1.05
    fig.add_trace(go.Scatter3d(
        x=[es_pos[0]], y=[es_pos[1]], z=[es_pos[2]], mode='text',
        text=['ES'], textfont=dict(size=15, color='yellow'), name='El Segundo'
    ))

    fig.update_layout(
        title=title,
        scene=dict(
            xaxis_title='X (m)', yaxis_title='Y (m)', zaxis_title='Z (m)',
            aspectmode='data'
        ),
        margin=dict(r=20, b=10, l=10, t=40),
        legend_title_text='Objects'
    )

    return fig
\end{verbatim}

\section{Code Listing: plotting\_vectors.py}
\begin{verbatim}
import numpy as np
from datetime import datetime
from typing import Dict, Any

import plotly.graph_objects as go

EARTH_RADIUS = 6378137.0

def plot_pointing_vectors(
    data_struct: Dict[str, Any],
    title: str,
    plot_time: datetime
) -> go.Figure:
    """
    Creates a 3D plot of satellites with pointing vectors.
    """
    sat_positions = data_struct['satellites']['position'].copy()
    sat_pointing = data_struct['satellites']['pointing'].copy()
    num_sats = data_struct['counts']['satellites']

    if num_sats == 0:
        print("No satellites to plot.")
        return go.Figure()

    fig = go.Figure()

    fig.add_trace(go.Scatter3d(
        x=sat_positions[:, 0], y=sat_positions[:, 1], z=sat_positions[:, 2],
        mode='markers',
        marker=dict(size=5, color='blue', opacity=0.8),
        name='Satellites'
    ))

    vector_scale = 0.5 * EARTH_RADIUS
    for i in range(num_sats):
        start_point = sat_positions[i]
        pointing_vec = sat_pointing[i]
        norm = np.linalg.norm(pointing_vec)
        unit_vec = pointing_vec / norm if norm > 0 else np.array([0, 0, 0])
        end_point = start_point + unit_vec * vector_scale
        fig.add_trace(go.Scatter3d(
            x=[start_point[0], end_point[0]],
            y=[start_point[1], end_point[1]],
            z=[start_point[2], end_point[2]],
            mode='lines',
            line=dict(color='red', width=3),
            showlegend=False
        ))

    u_sphere = np.linspace(0, 2 * np.pi, 100)
    v_sphere = np.linspace(0, np.pi, 100)
    x_earth = EARTH_RADIUS * np.outer(np.cos(u_sphere), np.sin(v_sphere))
    y_earth = EARTH_RADIUS * np.outer(np.sin(u_sphere), np.sin(v_sphere))
    z_earth = EARTH_RADIUS * np.outer(np.ones(np.size(u_sphere)), np.cos(v_sphere))
    fig.add_trace(go.Surface(
        x=x_earth, y=y_earth, z=z_earth,
        colorscale='Blues', showscale=False, opacity=0.5, name='Earth'
    ))

    fig.update_layout(
        title=title,
        scene=dict(
            xaxis_title='X (m)', yaxis_title='Y (m)', zaxis_title='Z (m)',
            aspectmode='data'
        ),
        margin=dict(r=20, b=10, l=10, t=40)
    )

    return fig
\end{verbatim}

\section{Code Listing: pointing.py}
\begin{verbatim}
import numpy as np
from typing import Dict, Any

from constants import POINTING_COUNT_IDX, POINTING_PLACE_IDX
from pointing_vectors import pointing_vectors

def pointing_place_update(data_struct: Dict[str, Any]) -> None:
    """
    Increments the pointing place for all satellites, wrapping around if necessary.
    Basically this usges data_struct['satellites']['pointing_state'] and wraps to
    0 if its POINTING_COUNT_IDX or greater.
    """

    pointing_state = data_struct['satellites']['pointing_state']
    pointing_counts = pointing_state[:, POINTING_COUNT_IDX]

    # Increment pointing place
    pointing_state[:, POINTING_PLACE_IDX] += 1

    # Wrap around where place >= count
    wrap_around_indices = np.where(pointing_state[:, POINTING_PLACE_IDX] >= pointing_counts)
    pointing_state[wrap_around_indices, POINTING_PLACE_IDX] = 0

def jerk(data_struct: Dict[str, Any], satellite_number: int) -> Dict[str, Any]:
    """
    Moves the pointing vector of a specific satellite by 0.3 radians in a
    random direction.

    This function applies a random rotation to the satellite's pointing vector
    using a simplified version of Rodrigues' rotation formula.

    Args:
        data_struct: The main simulation data dictionary.
        satellite_number: The index of the satellite to modify.

    Returns:
        The modified data_struct with the updated pointing vector.
    """
    p = data_struct['satellites']['pointing'][satellite_number]
    p_norm = p / np.linalg.norm(p)

    while True:
        r = np.random.rand(3) - 0.5
        k = np.cross(p_norm, r)
        norm_k = np.linalg.norm(k)
        if norm_k > 1e-9:
            k_hat = k / norm_k
            break

    theta = 0.3
    cos_theta = np.cos(theta)
    sin_theta = np.sin(theta)

    p_new = p_norm * cos_theta + np.cross(k_hat, p_norm) * sin_theta

    data_struct['satellites']['pointing'][satellite_number] = p_new / np.linalg.norm(p_new)

    return data_struct

def generate_pointing_sphere(data_struct: Dict[str, Any], n_points: int) -> None:
    """
    Generates a pointing sphere with n_points and stores it in the data_struct.
    The index is ['pointing_spheres'][n]
    If a sphere with the same number of points already exists, this function does nothing.
    """
    if n_points not in data_struct['pointing_spheres']:
        print(f"Generating pointing sphere with {n_points} points...")
        data_struct['pointing_spheres'][n_points] = pointing_vectors(n_points)

def update_satellite_pointing(data_struct: Dict[str, Any]) -> None:
    """
    Updates the pointing vector for each satellite based on its pointing state.
    """
    num_sats = data_struct['counts']['satellites']
    if num_sats == 0:
        return

    pointing_state = data_struct['satellites']['pointing_state']
    pointing_vectors_all = data_struct['satellites']['pointing']

    for i in range(num_sats):
        count = int(pointing_state[i, POINTING_COUNT_IDX])
        place = int(pointing_state[i, POINTING_PLACE_IDX])

        if count > 0:
            if count not in data_struct['pointing_spheres']:
                raise ValueError(f"Pointing sphere for {count} points not generated.")

            grid = data_struct['pointing_spheres'][count]
            pointing_vectors_all[i] = grid[place]

def find_and_jerk_blind_satellites(data_struct: Dict[str, Any]) -> Dict[str, Any]:
    """
    Finds satellites with no visibility and applies the 'jerk' function to them.

    This function iterates through the visibility table. If any satellite (column)
    has no visible fixed points (i.e., the column sum is 0), the `jerk`
    function is called to randomly adjust its pointing vector.

    Args:
        data_struct: The main simulation data dictionary.

    Returns:
        The modified data_struct.
    """
    visibility_table = data_struct['fixedpoints']['visibility']

    column_sums = np.sum(visibility_table, axis=0)
    blind_satellite_indices = np.where(column_sums == 0)[0]

    for sat_index in blind_satellite_indices:
        print(f"Satellite {sat_index} has no visible points. Applying jerk.")
        data_struct = jerk(data_struct, sat_index)

    return data_struct
\end{verbatim}

\section{Code Listing: pointing\_vectors.py}
\begin{verbatim}
import numpy as np
import plotly.graph_objects as go


def pointing_vectors(n: int) -> np.ndarray:
    """
    Generates n equally spaced points on a unit sphere using the Fibonacci lattice algorithm.

    Args:
        n: The number of points to generate.

    Returns:
        A NumPy array of shape (n, 3) for the Cartesian coordinates of the points.
    """
    if not isinstance(n, int) or n <= 0:
        raise ValueError("n must be a positive integer.")

    indices = np.arange(0, n, dtype=float) + 0.5
    z = 1 - 2 * indices / n
    radius_xy = np.sqrt(1 - z**2)
    golden_angle = np.pi * (3. - np.sqrt(5.))
    theta = golden_angle * indices
    x = radius_xy * np.cos(theta)
    y = radius_xy * np.sin(theta)

    unit_vectors = np.stack([x, y, z], axis=1)
    return unit_vectors

def plot_vectors_on_sphere(vectors: np.ndarray, title: str) -> go.Figure:
    """
    Creates a 3D plot of vectors on a sphere.

    Args:
        vectors: A NumPy array of shape (n, 3) representing the vectors.
        title: The title of the plot.

    Returns:
        A Plotly figure object.
    """
    if vectors.ndim != 2 or vectors.shape[1] != 3:
        raise ValueError("vectors array must have shape (n, 3).")

    fig = go.Figure()

    fig.add_trace(go.Scatter3d(
        x=vectors[:, 0], y=vectors[:, 1], z=vectors[:, 2],
        mode='markers',
        marker=dict(
            size=2,
            color='red',
            opacity=0.8
        ),
        name='Vectors'
    ))

    u_sphere = np.linspace(0, 2 * np.pi, 100)
    v_sphere = np.linspace(0, np.pi, 100)
    x_earth = 1.0 * np.outer(np.cos(u_sphere), np.sin(v_sphere))
    y_earth = 1.0 * np.outer(np.sin(u_sphere), np.sin(v_sphere))
    z_earth = 1.0 * np.outer(np.ones(np.size(u_sphere)), np.cos(v_sphere))
    fig.add_trace(go.Surface(
        x=x_earth, y=y_earth, z=z_earth,
        colorscale='Blues', showscale=False, opacity=0.5, name='Sphere'
    ))

    fig.update_layout(
        title=title,
        scene=dict(
            xaxis_title='X', yaxis_title='Y', zaxis_title='Z',
            aspectmode='data'
        ),
        margin=dict(r=20, b=10, l=10, t=40),
        legend_title_text='Objects'
    )

    return fig
\end{verbatim}

\section{Code Listing: propagation.py}
\begin{verbatim}
import numpy as np
from datetime import datetime, timezone
from typing import Dict, Any, List, Tuple
from sgp4.api import Satrec
from astropy.time import Time
from astropy.coordinates import get_body, GCRS
import astropy.units as u

from constants import (
    ORBITAL_A_IDX, ORBITAL_E_IDX, ORBITAL_I_IDX,
    ORBITAL_RAAN_IDX, ORBITAL_ARGP_IDX, ORBITAL_M_IDX
)

def add_satellites_from_tle(sim_data: Dict[str, Any], tle_file_path: str, sat_category: str) -> None:
    """
    Adds and initializes a category of satellites from a TLE file.
    Mote although the TLEs are loaded, positions etc. are not.

    Args:
        sim_data: The main simulation data dictionary.
        tle_file_path: Path to the TLE file.
        sat_category: The key for this satellite category (e.g., 'satellites').

    Data added to the set_category element of sim_data
    position
    velocity
    acceleration
    orbital_elements
    epochs
    pointing
    """
    orbital_elements, epochs = readtle(tle_file_path)
    num_sats = len(epochs)

    sim_data['counts'][sat_category] = num_sats
    sim_data[sat_category] = {
        'position': np.zeros((num_sats, 3), dtype=float),
        'velocity': np.zeros((num_sats, 3), dtype=float),
        'acceleration': np.zeros((num_sats, 3), dtype=float),
        'orbital_elements': orbital_elements,
        'epochs': epochs,
        'pointing': np.zeros((num_sats, 3), dtype=float),
        'pointing_state': np.zeros((num_sats, 2), dtype=int),
        'detector': np.zeros((num_sats, 7), dtype=float),
    }

def celestial_update(data_struct: Dict[str, Any], time_date: datetime) -> Dict[str, Any]:
    """
    Calculates and updates the positions of the Sun and Moon.
    """
    if time_date.tzinfo is None:
        raise ValueError("time_date must be timezone-aware.")

    astro_time = Time(time_date)
    sun_coords = get_body("sun", astro_time)
    sun_gcrs = sun_coords.transform_to(GCRS(obstime=astro_time))
    moon_coords = get_body("moon", astro_time)
    moon_gcrs = moon_coords.transform_to(GCRS(obstime=astro_time))

    celestial_pos = data_struct['celestial']['position']
    celestial_pos[0] = sun_gcrs.cartesian.xyz.to(u.m).value
    celestial_pos[1] = moon_gcrs.cartesian.xyz.to(u.m).value

    return data_struct

def readtle(tle_file_path: str) -> Tuple[np.ndarray, List[datetime]]:
    """
    Reads a TLE file and extracts orbital elements and epochs for each satellite.

    The array returned ahd the orbital elements in "canonical" order.
    """
    orbital_elements_list = []
    epochs_list = []
    with open(tle_file_path, 'r') as f:
        lines = f.readlines()

    for i in range(0, len(lines), 3):
        line1 = lines[i+1].strip()
        line2 = lines[i+2].strip()

        satellite = Satrec.twoline2rv(line1, line2)

        jd, fr = satellite.jdsatepoch, satellite.jdsatepochF
        epoch_dt = Time(jd, fr, format='jd', scale='utc').to_datetime(timezone.utc)
        epochs_list.append(epoch_dt)

        a = satellite.a * satellite.radiusearthkm * 1000.0
        e = satellite.ecco
        inc = satellite.inclo
        raan = satellite.nodeo
        argp = satellite.argpo
        M = satellite.mo

        elements = np.zeros(6)
        elements[ORBITAL_A_IDX] = a
        elements[ORBITAL_E_IDX] = e
        elements[ORBITAL_I_IDX] = inc
        elements[ORBITAL_RAAN_IDX] = raan
        elements[ORBITAL_ARGP_IDX] = argp
        elements[ORBITAL_M_IDX] = M
        orbital_elements_list.append(elements)

    return np.array(orbital_elements_list, dtype=float), epochs_list

def propagate_satellites_new(data_struct: Dict[str, Any], time_date: datetime,
                            sat_category: str = None) -> Dict[str, Any]:
    """
    Updates satellite positions and pointing vectors based on their orbital elements.
    """
    MU_EARTH = 3.986004418e14
    time_date_timestamp = time_date.timestamp()

    if sat_category:
        categories = [sat_category]
    else:
        categories = ['satellites', 'red_satellites']

    for category in categories:
        if category not in data_struct['counts'] or data_struct['counts'][category] == 0:
            continue

        elements = data_struct[category]['orbital_elements']
        epochs = data_struct[category]['epochs']

        epoch_timestamps = np.array([e.timestamp() for e in epochs])
        delta_t_array = time_date_timestamp - epoch_timestamps

        a = elements[:, ORBITAL_A_IDX]
        e = elements[:, ORBITAL_E_IDX]
        i = elements[:, ORBITAL_I_IDX]
        raan = elements[:, ORBITAL_RAAN_IDX]
        argp = elements[:, ORBITAL_ARGP_IDX]
        M0 = elements[:, ORBITAL_M_IDX]

        n = np.sqrt(MU_EARTH / a**3)
        M = (M0 + n * delta_t_array) % (2 * np.pi)

        E = M.copy()
        for _ in range(10):
            f_E = E - e * np.sin(E) - M
            f_prime_E = 1 - e * np.cos(E)
            f_prime_E[f_prime_E == 0] = 1e-10
            E = E - f_E / f_prime_E

        tan_nu_half = np.sqrt((1 + e) / (1 - e)) * np.tan(E / 2)
        nu = 2 * np.arctan(tan_nu_half)

        r = a * (1 - e * np.cos(E))

        x_pqw = r * np.cos(nu)
        y_pqw = r * np.sin(nu)

        cos_raan = np.cos(raan)
        sin_raan = np.sin(raan)
        cos_argp = np.cos(argp)
        sin_argp = np.sin(argp)
        cos_i = np.cos(i)
        sin_i = np.sin(i)

        P_x = cos_argp * cos_raan - sin_argp * sin_raan * cos_i
        P_y = cos_argp * sin_raan + sin_argp * cos_raan * cos_i
        P_z = sin_argp * sin_i

        Q_x = -sin_argp * cos_raan - cos_argp * sin_raan * cos_i
        Q_y = -sin_argp * sin_raan + cos_argp * cos_raan * cos_i
        Q_z = cos_argp * sin_i

        x_gcrs = x_pqw * P_x + y_pqw * Q_x
        y_gcrs = x_pqw * P_y + y_pqw * Q_y
        z_gcrs = x_pqw * P_z + y_pqw * Q_z

        positions = np.vstack((x_gcrs, y_gcrs, z_gcrs)).T
        data_struct[category]['position'] = positions

        norms = np.linalg.norm(positions, axis=1)[:, np.newaxis]
        norms[norms == 0] = 1.0
        data_struct[category]['pointing'] = positions / norms

    return data_struct
\end{verbatim}

\section{Code Listing: radiometry\_calcs.py}
\begin{verbatim}
import math
import scipy.integrate
import scipy.constants as const
import numpy as np
import plotly.graph_objects as go

def mag(x: float) -> float:
    """
    Calculates a magnitude value from a linear ratio.
    Uses the formula: magnitude = -2.5 * log10(ratio)
    """
    if x <= 0:
        raise ValueError("Input for mag() must be positive.")
    return -2.5 * math.log10(x)

def amag(x: float) -> float:
    """
    Calculates the linear ratio from a magnitude value.
    This is the inverse of the mag() function.
    Uses the formula: ratio = 10**(-0.4 * magnitude)
    """
    return 10**(-0.4 * x)

def _planck_law(wav_m: float, temp_k: float) -> float:
    """
    Helper function for Planck's law for spectral radiance.
    Args:
        wav_m: Wavelength in meters.
        temp_k: Temperature in Kelvin.
    Returns:
        Spectral radiance in W / (m^2 * sr * m).
    """
    if wav_m <= 0:
        return 0.0
    exponent = (const.h * const.c) / (wav_m * const.k * temp_k)
    if exponent > 700:
        return 0.0
    numerator = 2.0 * const.h * const.c**2
    denominator = (wav_m**5) * (math.expm1(exponent))
    if denominator == 0:
        return 0.0
    return numerator / denominator

def blackbody_flux(
    temperature: float,
    lambda_short: float,
    lambda_long: float
) -> float:
    """
    Numerically computes the integrated spectral radiance of
    a blackbody over a given wavelength band.

    Requires the scipy library.

    Args:
        temperature: The temperature of the blackbody in Kelvin.
        lambda_short: The short wavelength of the band in microns.
        lambda_long: The long wavelength of the band in microns.

    Returns:
        The integrated spectral radiance in units of
        Watts / (m^2 * steradian).
    """
    if lambda_short >= lambda_long:
        raise ValueError("lambda_short must be less than lambda_long.")
    lambda_short_m = lambda_short * 1e-6
    lambda_long_m = lambda_long * 1e-6
    integrated_radiance, _ = scipy.integrate.quad(
        lambda wav: _planck_law(wav, temperature),
        lambda_short_m,
        lambda_long_m
    )
    return integrated_radiance

def stefan_boltzmann_law(temperature: float) -> float:
    """
    Calculates the total power radiated per unit area by a
    blackbody using the Stefan-Boltzmann law.

    Args:
        temperature: The temperature of the blackbody in Kelvin.

    Returns:
        The total radiated power per unit area in W / m^2.
    """
    return const.sigma * (temperature**4)

def plot_blackbody_spectrum(temperature: float):
    """
    Plots the spectral radiance of a blackbody from
    0.5 to 30 microns.

    Args:
        temperature: The temperature of the blackbody in Kelvin.
    """
    wavelengths_microns = np.geomspace(0.5, 30, 500)
    wavelengths_m = wavelengths_microns * 1e-6
    spectral_radiance = [_planck_law(w, temperature) * 1e-6 for w in wavelengths_m]
    fig = go.Figure()
    fig.add_trace(go.Scatter(
        x=wavelengths_microns,
        y=spectral_radiance,
        mode='lines',
        name=f'{temperature} K Blackbody'
    ))
    fig.update_layout(
        title_text=f'Blackbody Spectrum ({temperature} K)',
        xaxis_title="Wavelength (microns)",
        yaxis_title="Spectral Radiance (W / m^2 / sr / micron)",
        xaxis_type="log",
        yaxis_type="log",
        template="plotly_white"
    )
    fig.show()

def plot_blackbody_spectrum_visible_nir(temperature: float):
    """
    Plots the spectral radiance of a blackbody from
    0.1 to 1 micron.

    Args:
        temperature: The temperature of the blackbody in Kelvin.
    """
    wavelengths_microns = np.linspace(0.1, 1.0, 500)
    wavelengths_m = wavelengths_microns * 1e-6
    spectral_radiance = [_planck_law(w, temperature) * 1e-6 for w in wavelengths_m]
    fig = go.Figure()
    fig.add_trace(go.Scatter(
        x=wavelengths_microns,
        y=spectral_radiance,
        mode='lines',
        name=f'{temperature} K Blackbody'
    ))
    fig.update_layout(
        title_text=f'Blackbody Spectrum ({temperature} K) - Visible/NIR',
        xaxis_title="Wavelength (microns)",
        yaxis_title="Spectral Radiance (W / m^2 / sr / micron)",
        template="plotly_white"
    )
    fig.show()
\end{verbatim}

\section{Code Listing: radiometry\_data.py}
\begin{verbatim}
import scipy.constants as const

# --- Physical Constants ---
AU_M = const.au  # Astronomical Unit in meters
RSUN_M = 6.957e8 # Radius of the Sun in meters

# --- Radiometric Data ---

# Data is compiled for standard astronomical filters.
#
# Magnitudes are apparent magnitudes in the AB or Vega system.
# Sky brightness is for a dark site at new moon, in
# magnitudes per square arcsecond. For space-based IR
# telescopes, this is the zodiacal + telescope background.
# For ground-based IR, it's dominated by thermal emission.
# Wavelengths and bandwidths are in nanometers (nm).
# Zero point is the photon flux for a 0-magnitude
# object in photons per second per square meter.

FILTER_DATA = {
    # Johnson-Cousins UBVRI Filters
    'U': {
        'sun': -26.03,
        'sky': 22.0,
        'space': 22.9,
        'central_wavelength': 365.0,
        'bandwidth': 66.0,
        'zero_point': 4.96e9,
    },
    'B': {
        'sun': -26.13,
        'sky': 22.7,
        'space': 23.3,
        'central_wavelength': 445.0,
        'bandwidth': 94.0,
        'zero_point': 1.36e10,
    },
    'V': {
        'sun': -26.78,
        'sky': 21.9,
        'space': 23,
        'central_wavelength': 551.0,
        'bandwidth': 88.0,
        'zero_point': 8.79e9,
    },
    'R': {
        'sun': -27.11,
        'sky': 21.0,
        'space': 22.0,
        'central_wavelength': 658.0,
        'bandwidth': 138.0,
        'zero_point': 9.74e9,
    },
    'I': {
        'sun': -27.47,
        'sky': 20.0,
        'space': 21.2,
        'central_wavelength': 806.0,
        'bandwidth': 149.0,
        'zero_point': 7.11e9,
    },
    # Near-Infrared JHK Filters
    'J': {
        'sun': -27.91,
        'sky': 16.0,
        'space': 21.5,
        'central_wavelength': 1235.0,
        'bandwidth': 162.0,
        'zero_point': 3.15e9,
    },
    'H': {
        'sun': -28.25,
        'sky': 14.0,
        'space': 20.5,
        'central_wavelength': 1646.0,
        'bandwidth': 251.0,
        'zero_point': 2.35e9,
    },
    'K': {
        'sun': -28.29,
        'sky': 13.0,
        'space': 19.5,
        'central_wavelength': 2159.0,
        'bandwidth': 262.0,
        'zero_point': 1.17e9,
    },
    # SDSS Filters
    'g': {
        'sun': -26.47,
        'sky': 21.8,
        'space': 23.0,
        'central_wavelength': 477.0,
        'bandwidth': 137.9,
        'zero_point': 1.582e10,
    },
    'r': {
        'sun': -26.93,
        'sky': 20.8,
        'space': 22.2,
        'central_wavelength': 623.1,
        'bandwidth': 138.2,
        'zero_point': 1.214e10,
    },
    'i': {
        'sun': -27.05,
        'sky': 20.2,
        'space': 21.5,
        'central_wavelength': 762.5,
        'bandwidth': 153.5,
        'zero_point': 1.103e10,
    },
    'z': {
        'sun': -27.07,
        'sky': 19.0,
        'space': 20.8,
        'central_wavelength': 913.4,
        'bandwidth': 140.9,
        'zero_point': 8.455e9,
    },
    # Ground-Based Mid-IR Filters
    'L': {
        'sun': None,
        'sky': 3.5,
        'central_wavelength': 3500.0,
        'bandwidth': 600.0,
        'zero_point': 9.4e9,
    },
    'M': {
        'sun': None,
        'sky': 0.0,
        'central_wavelength': 4800.0,
        'bandwidth': 600.0,
        'zero_point': 6.8e9,
    },
    'N': {
        'sun': None,
        'sky': -6.0,
        'central_wavelength': 10200.0,
        'bandwidth': 5000.0,
        'zero_point': 2.7e10,
    },
    # JWST MIRI Filters
    'F560W': { 'sun': None,
               'sky': 24.4,
               'central_wavelength': 5600.0,
               'bandwidth': 1000.0,
               'zero_point': 9.78e9 },
    'F770W': { 'sun': None,
               'sky': 24.4,
               'central_wavelength': 7700.0,
               'bandwidth': 1950.0, '
               zero_point': 1.39e10 },
    'F1000W': { 'sun': None,
                'sky': 23.6,
                'central_wavelength': 10000.0,
                'bandwidth': 1800.0, '
                zero_point': 9.86e9 },
    'F1500W': { 'sun': None,
                'sky': 21.9,
                'central_wavelength': 15000.0,
                'bandwidth': 2940.0, '
                zero_point': 1.08e10 },
    'F2550W': { 'sun': None,
                'sky': 19.4,
                'central_wavelength': 25500.0,
                'bandwidth': 4050.0,
                'zero_point': 8.70e9 }
}
\end{verbatim}

\section{Code Listing: simulation.py}
\begin{verbatim}
import numpy as np
from datetime import datetime, timezone
from typing import Dict, Any
from generate_log_spherical_points import generate_log_spherical_points
from constants import *

def create_empty_simulation(start_time: datetime, delta_time: float = 60.0) -> Dict[str, Any]:
    """
    Initializes a minimal, empty data structure for a space simulation.

    Args:
        start_time: The starting time and date of the simulation. This must be a
                    timezone-aware datetime object set to UTC.
        delta_time: The time step for the simulation in seconds.

    Returns:
        A dictionary representing the basic simulation state.
    """
    if not isinstance(start_time, datetime):
        raise TypeError("start_time must be a datetime object.")
    if start_time.tzinfo is None:
        raise ValueError("start_time must be timezone-aware. Please set tzinfo.")

    simulation_data: Dict[str, Any] = {
        'start_time': start_time,
        'delta_time': delta_time,
        'counts': {},
        'pointing_spheres': {},
    }
    return simulation_data

def add_celestial_bodies(sim_data: Dict[str, Any]) -> None:
    """
    Adds celestial body structures (for Sun and Moon) to the simulation data.

    Args:
        sim_data: The simulation data dictionary.
    """
    sim_data['counts']['celestial'] = 2  # Sun and Moon
    sim_data['celestial'] = {
        'position': np.zeros((2, 3), dtype=float),
        'velocity': np.zeros((2, 3), dtype=float),
        'acceleration': np.zeros((2, 3), dtype=float),
    }

def add_fixed_points(sim_data: Dict[str, Any], num_points: int = 100) -> None:
    """
    Adds a structure for fixed reference points in the GCRS frame.

    Args:
        sim_data: The simulation data dictionary.
        num_points: The number of fixed points to generate.
    """
    sim_data['counts']['fixedpoints'] = num_points
    sim_data['fixedpoints'] = {
        'position': generate_log_spherical_points(
            num_points=num_points,
            inner_radius=2000000,
            outer_radius=84328000
        )[0],
        'visibility': np.zeros((num_points, 0), dtype=int) # Visibility will be resized later
    }
\end{verbatim}

\section{Code Listing: visibility.py}
\begin{verbatim}
import numpy as np
from typing import Dict, Any, Tuple, Optional

from constants import (
    EARTH_RADIUS, MOON_RADIUS, DETECTOR_SOLAR_EXCL_IDX,
    DETECTOR_LUNAR_EXCL_IDX, DETECTOR_EARTH_EXCL_IDX
)

def solarexclusion(data_struct: Dict[str, Any]) -> Tuple[np.ndarray, np.ndarray]:
    """
    Calculates solar exclusion for all satellites based on their pointing vectors.

    This function operates in a vectorized manner on all satellites in the
    'satellites' category. It computes the angle between each satellite's
    pointing vector and the vector from the satellite to the Sun.

    Args:
        data_struct: The main simulation data dictionary.

    Returns:
        A tuple containing:
        - exclusion_vector (np.ndarray): An array of the same length as the
          number of satellites. An element is 1 if the satellite is within
          the solar exclusion angle, 0 otherwise.
        - angle_vector (np.ndarray): An array containing the calculated angle
          in radians for each satellite.
    """
    num_sats = data_struct['counts']['satellites']
    if num_sats == 0:
        return np.array([]), np.array([])

    sun_pos = data_struct['celestial']['position'][0]
    sat_pos = data_struct['satellites']['position']
    sat_pointing = data_struct['satellites']['pointing']
    solar_exclusion_angles = data_struct['satellites']['detector'][:, DETECTOR_SOLAR_EXCL_IDX]

    vec_sat_to_sun = sun_pos - sat_pos

    norm_sat_to_sun = np.linalg.norm(vec_sat_to_sun, axis=1)
    norm_sat_pointing = np.linalg.norm(sat_pointing, axis=1)

    valid_norms = (norm_sat_to_sun > 1e-9) & (norm_sat_pointing > 1e-9)

    angle_vector = np.full(num_sats, np.pi)

    if np.any(valid_norms):
        dot_product = np.einsum('ij,ij->i', vec_sat_to_sun[valid_norms],
                                sat_pointing[valid_norms])
        cos_angle = dot_product / (norm_sat_to_sun[valid_norms] *
                                   norm_sat_pointing[valid_norms])
        cos_angle = np.clip(cos_angle, -1.0, 1.0)
        angle_vector[valid_norms] = np.arccos(cos_angle)

    exclusion_vector = (angle_vector < solar_exclusion_angles).astype(int)

    return exclusion_vector, angle_vector

def exclusion(
    data_struct: Dict[str, Any],
    satellite_index: int,
    print_debug: bool = False
) -> int:
    """
    Determines if a satellite's pointing vector is excluded by the Sun, Moon, or Earth.

    Args:
        data_struct: The main simulation data dictionary.
        satellite_index: The index of the satellite to check.
        print_debug: If True, prints detailed debug information for the calculation.

    Returns:
        0 if the satellite's view is excluded, 1 otherwise.
    """
    sat_pos = data_struct['satellites']['position'][satellite_index]
    sat_pointing = data_struct['satellites']['pointing'][satellite_index]
    sun_pos = data_struct['celestial']['position'][0]
    moon_pos = data_struct['celestial']['position'][1]

    detector_props = data_struct['satellites']['detector'][satellite_index]
    solar_excl_angle = detector_props[DETECTOR_SOLAR_EXCL_IDX]
    lunar_excl_angle = detector_props[DETECTOR_LUNAR_EXCL_IDX]
    earth_excl_angle = detector_props[DETECTOR_EARTH_EXCL_IDX]

    vec_to_sun = sun_pos - sat_pos
    vec_to_moon = moon_pos - sat_pos
    vec_to_earth = -sat_pos

    dist_to_sun = np.linalg.norm(vec_to_sun)
    dist_to_moon = np.linalg.norm(vec_to_moon)
    dist_to_earth = np.linalg.norm(vec_to_earth)

    u_vec_to_sun = vec_to_sun / dist_to_sun if dist_to_sun > 0 else np.array([0.,0.,0.])
    u_vec_to_moon = vec_to_moon / dist_to_moon if dist_to_moon > 0 else np.array([0.,0.,0.])
    u_vec_to_earth = vec_to_earth / dist_to_earth if dist_to_earth > 0 else np.array([0.,0.,0.])

    norm_pointing = np.linalg.norm(sat_pointing)
    u_sat_pointing = sat_pointing / norm_pointing if norm_pointing > 0 else np.array([0.,0.,0.])

    sun_flag, moon_flag, earth_flag = False, False, False

    cos_angle_sun = np.clip(np.dot(u_sat_pointing, u_vec_to_sun), -1.0, 1.0)
    angle_sun = np.arccos(cos_angle_sun)
    if angle_sun < solar_excl_angle:
        sun_flag = True

    cos_angle_moon = np.clip(np.dot(u_sat_pointing, u_vec_to_moon), -1.0, 1.0)
    angle_moon = np.arccos(cos_angle_moon)
    apparent_radius_moon = np.arctan(MOON_RADIUS / dist_to_moon) if dist_to_moon > 0 else 0
    if (angle_moon - apparent_radius_moon) < lunar_excl_angle:
        moon_flag = True

    cos_angle_earth = np.clip(np.dot(u_sat_pointing, u_vec_to_earth), -1.0, 1.0)
    angle_earth = np.arccos(cos_angle_earth)
    apparent_radius_earth = np.arctan(EARTH_RADIUS / dist_to_earth) if dist_to_earth > 0 else 0
    if (angle_earth - apparent_radius_earth) < earth_excl_angle:
        earth_flag = True

    if print_debug:
        print(f"--- Exclusion Debug for Satellite {satellite_index} ---")
        print(f"  - Sun Flag:   {sun_flag}, Angle: {np.rad2deg(angle_sun):.2f} deg, "
              f"Excl: {np.rad2deg(solar_excl_angle):.2f} deg")
        print(f"  - Moon Flag:  {moon_flag}, Angle: {np.rad2deg(angle_moon):.2f} deg, "
              f"Excl: {np.rad2deg(lunar_excl_angle):.2f} deg")
        print(f"  - Earth Flag: {earth_flag}, Angle: {np.rad2deg(angle_earth):.2f} deg, "
              f"Excl: {np.rad2deg(earth_excl_angle):.2f} deg")

    is_excluded = sun_flag or moon_flag or earth_flag
    return 0 if is_excluded else 1

def update_visibility_table(
    data_struct: Dict[str, Any],
    print_debug_for_sat: Optional[int] = None
) -> None:
    """
    Updates the visibility table for each satellite against each fixed point.

    Args:
        data_struct: The main simulation data dictionary.
        print_debug_for_sat: If an integer is provided, the `exclusion` function's
                             debug printout will be enabled for that satellite index.
    """
    num_satellites = data_struct['counts']['satellites']
    fixed_points = data_struct['fixedpoints']['position']
    num_fixed_points = len(fixed_points)
    visibility_table = data_struct['fixedpoints']['visibility']

    # Ensure the visibility table has the correct shape
    if visibility_table.shape != (num_fixed_points, num_satellites):
        data_struct['fixedpoints']['visibility'] = np.zeros((num_fixed_points, num_satellites),
                                                            dtype=int)
        visibility_table = data_struct['fixedpoints']['visibility']

    if num_satellites == 0 or num_fixed_points == 0:
        return

    satellite_positions = data_struct['satellites']['position']

    for i in range(num_satellites):
        sat_pos = satellite_positions[i]
        should_print_debug = (i == print_debug_for_sat)
        for j in range(num_fixed_points):
            fixed_point_pos = fixed_points[j]
            pointing_vector = fixed_point_pos - sat_pos
            data_struct['satellites']['pointing'][i] = pointing_vector
            visibility_table[j, i] = exclusion(data_struct, i,
                                               print_debug=should_print_debug)
\end{verbatim}

\end{document}